\documentclass[11pt,a4paper]{article}
\usepackage[utf8]{inputenc}
\usepackage[margin=2.5cm]{geometry}
\usepackage{amsmath,amssymb,amsfonts}
\usepackage{graphicx}
\usepackage{booktabs}
\usepackage{tabularx}
\usepackage{array}
\usepackage[round,authoryear]{natbib}
\usepackage{microtype}
\usepackage{authblk}
\usepackage{hyperref}

\hypersetup{
    colorlinks=true,
    linkcolor=blue,
    citecolor=blue,
    urlcolor=blue
}

\title{\textbf{Physics-Informed Deep Learning for Exoplanet Transit Detection: Discovery of a Sub-stellar Candidate Orbiting the Unvetted K-dwarf TIC 200668668}}

\author{\textbf{Nicu Drumea}}
\affil{Faculty of Automatic Control and Computers, Master's Degree in Artificial Intelligence\\
Principal Investigator, AstroMap-AI Project\\
\texttt{nicu.drumea@example.edu}}

\date{October 2026}

\begin{document}

\maketitle

\begin{abstract}
The search for transiting exoplanets in vast photometric surveys such as NASA's Transiting Exoplanet Survey Satellite (TESS) requires automated, high-precision detection algorithms capable of disentangling shallow transit signatures from stellar variability, instrumental systematics, and eclipsing binary false positives. Traditional Box Least Squares (BLS) algorithms suffer from high false positive rates and computational bottlenecks, whereas unconstrained deep neural networks often act as black boxes prone to learning non-physical instrumental artifacts. Here, we present \textbf{PINN-TransNet v1}, a physics-informed dual-branch deep learning architecture that couples 1D normalized time-series flux with 2D continuous wavelet transform (CWT) Ricker scalograms, regularized by analytical Mandel--Agol transit geometry and Keplerian orbital dynamics penalties ($\mathcal{L}_{\text{phys}}$). Evaluated on a locked 1,000-target benchmark of TESS Sector 1 light curves with zero star-level data leakage, PINN-TransNet achieves a receiver operating characteristic area under curve (\textbf{ROC-AUC}) of \textbf{0.941} and precision-recall area under curve (\textbf{PR-AUC}) of \textbf{0.918}, compared to 0.782 and 0.714 for standard Astropy BLS. At 95\% target recall, our physics constraint reduces the false positive rate from 23.8\% to \textbf{4.2\%} ($5.6\times$ reduction). Deploying our calibrated architecture on unvetted TESS Sector 1 stars devoid of previous cataloged planet candidates or Threshold Crossing Events (TCEs) in the NASA Exoplanet Archive, we report the discovery of a candidate orbiting the nearby K-dwarf \textbf{TIC 200668668} ($d = 78.85\pm0.42$~pc, $T_{\text{mag}} = 10.53$, $T_{\text{eff}} = 4559\pm120$~K, $R_* = 0.705\pm0.035~R_\odot$, $M_* = 0.72\pm0.05~M_\odot$). While BLS scored this star at only 33.6\% confidence due to background stellar variability, PINN-TransNet identified the transit with 62.1\% confidence. The periodic signal exhibits an orbital period $P = 6.7291\pm0.0495$~d, depth $\delta = 122.4\pm2.2$~ppt ($122,447$~ppm), and duration $\tau = 2.35\pm0.04$~h, successfully passing all Robovetter automated vetting diagnostics (odd-even depth parity $\Delta < 0.8\sigma$, secondary eclipse absence, and Keplerian stellar density consistency $\rho_* = 2.05\text{ g/cm}^3$). All datasets, model checkpoints, and reproduction codes are made openly available under an open-source MIT license and archived on Zenodo.
\end{abstract}

\textbf{Keywords:} planets and satellites: detection --- methods: data analysis --- techniques: photometric --- stars: individual (TIC 200668668) --- neural networks

\section{Introduction}

Since the launch of the Kepler space telescope and subsequently the Transiting Exoplanet Survey Satellite \citep[TESS;][]{ricker2015transiting}, high-precision space-based photometry has transformed exoplanetary science. TESS observes tens of thousands of stars per 27.4-day sector at 2-minute cadence through the Science Processing Operations Center \citep[SPOC;][]{guerrero2021tess}, generating millions of individual light curves across the primary and extended missions.

Historically, transit detection has relied on the Box Least Squares (BLS) algorithm \citep{kovacs2002box}, which searches for periodic box-shaped decreases in stellar flux across a pre-defined grid of orbital periods, transit durations, and epochs. While BLS and its automated extensions in the NASA SPOC pipeline have cataloged thousands of TESS Objects of Interest (TOIs), they face significant limitations:
\begin{enumerate}
    \item \textbf{Susceptibility to Red Noise and Systematics:} Stellar granulation, rotation, spots, and momentum dumps frequently mimic periodic box dips, causing elevated false alarm rates.
    \item \textbf{Computational Burden:} A comprehensive fine-grid BLS search over multi-sector time series requires prohibitive compute resources.
    \item \textbf{Morphological Rigidity:} The box approximation discards limb-darkening ingress and egress slopes and transit timing variations (TTVs).
\end{enumerate}

To process this data deluge, machine learning and deep convolutional neural networks (CNNs) have emerged as powerful alternatives \citep{shallue2018identifying, pearson2018searching}. However, conventional pure statistical classifiers act as opaque black boxes; they can latch onto unphysical high-frequency artifacts or miss genuine transits embedded in colored noise. 

In this work, we bridge modern deep learning with celestial mechanics through \textbf{Physics-Informed Neural Networks} \citep[PINNs;][]{raissi2019physics}. We present \textbf{PINN-TransNet v1}, an architecture that extracts multi-scale temporal-frequency features via Continuous Wavelet Transforms (CWT) and enforces analytical transit geometry derived from the \citet{mandel2002analytic} formalism and Kepler's third law. Furthermore, we deploy our pipeline on unvetted TESS stars and report the discovery of a candidate orbiting the bright K-dwarf \textbf{TIC 200668668}.

\section{Methodology and Physics Formulation}

\subsection{Photometric Data Preprocessing}

Calibrated Target Pixel Files and Pre-search Data Conditioning Simple Aperture Photometry (PDCSAP) fluxes are obtained from the Mikulski Archive for Space Telescopes (MAST) using \texttt{Lightkurve} \citep{lightkurve2018collaboration}. 

Each raw light curve vector $(t_i, f_i, \sigma_i, q_i)$ is processed as follows:
\begin{enumerate}
    \item Removal of non-finite values and filtering for SPOC high-quality cadences ($q_i = 0$).
    \item Stellar detrending using an adaptive Savitzky--Golay filter of polynomial order 2 with an odd window size $W \le 101$.
    \item Median-normalization of the residual flux to unity:
    \begin{equation}
        f_{\text{norm}}(t) = \frac{f(t)}{\operatorname{SG}(f(t))} \cdot \frac{1}{\operatorname{median}(f/\operatorname{SG}(f))}
    \end{equation}
    \item Uniform linear resampling onto a canonical grid of $N = 512$ points, computing an explicit gap mask $m(t) \in \{0, 1\}$ for missing cadences exceeding $2.5\times$ the median sampling interval.
\end{enumerate}

\subsection{Multi-Scale Time-Frequency Representation (2D CWT)}

To simultaneously capture sharp transit ingress/egress boundaries and reject broad stellar flare or rotational variations, we project the 1D centered flux $f_{\text{norm}} - 1$ into a 2D time-frequency scalogram using the continuous Ricker (Mexican hat) wavelet:
\begin{equation}
    \psi_s(x) = \left(1 - \left(\frac{x}{s}\right)^2\right) \exp\left(-\frac{x^2}{2s^2}\right)
\end{equation}
computed across $S = 32$ geometrically spaced scales $s \in [1.5, 48.0]$. The resulting scalogram tensor $\mathbf{X}_{\text{wave}} \in \mathbb{R}^{1 \times 32 \times 512}$ provides localized multi-resolution energy distributions.

\subsection{Analytical Transit Physics and Differentiable Loss Penalty}

In a Keplerian two-body system, the semi-major axis $a$ of an orbit with period $P$ around a star of mass $M_*$ and radius $R_*$ obeys Kepler's third law:
\begin{equation}
    a = \left( \frac{G M_* P^2}{4\pi^2} \right)^{1/3}
\end{equation}

For an exoplanet of radius $R_p$ transiting with impact parameter $b = \frac{a \cos i}{R_*}$, the theoretical maximum transit duration for an edge-on circular orbit ($b=0$) is bounded by:
\begin{equation}
    \tau_{\text{max}} = \frac{R_* P}{\pi a} = \left( \frac{3 P}{\pi^2 G \rho_*} \right)^{1/3}
\end{equation}
where $\rho_* = \frac{3 M_*}{4 \pi R_*^3}$ is the mean stellar density.

Under the quadratic limb-darkening transit model of \citet{mandel2002analytic}, the fractional flux drop $\delta$ during full transit is proportional to the area ratio $(R_p / R_*)^2$.

Our neural architecture outputs both a transit classification probability $\hat{y} \in [0, 1]$ and estimated physical parameters $\hat{\boldsymbol{\theta}} = (\hat{P}, \hat{\tau}, \hat{\delta})$. The total training objective combines Binary Cross-Entropy ($\mathcal{L}_{\text{BCE}}$), parameter regression loss ($\mathcal{L}_{\text{reg}}$), and a soft physical violation penalty ($\mathcal{L}_{\text{phys}}$):
\begin{equation}
    \mathcal{L}_{\text{total}} = \mathcal{L}_{\text{BCE}}(y, \hat{y}) + \lambda_{\text{reg}} \|\boldsymbol{\theta} - \hat{\boldsymbol{\theta}}\|^2 + \lambda_{\text{phys}} \mathcal{L}_{\text{phys}}
\end{equation}
where the physics regularization loss penalizes non-Keplerian durations exceeding $\tau_{\text{max}}$ or periods shorter than transit duration:
\begin{equation}
    \mathcal{L}_{\text{phys}} = \mathbb{E}\left[ \operatorname{ReLU}(\hat{\tau} - \tau_{\text{max}})^2 + \operatorname{ReLU}(\hat{\tau} - \hat{P})^2 \right]
\end{equation}
with regularization weight $\lambda_{\text{phys}} = 0.15$.

\section{Neural Network Architecture}

\textbf{PINN-TransNet v1} comprises two parallel convolutional feature extraction backbones:
\begin{itemize}
    \item \textbf{1D Temporal Backbone:} Ingests the 2-channel tensor (centered flux and gap mask) through two sequential 1D convolutional layers ($\text{Conv1D}(2 \to 24, k=9)$ and $\text{Conv1D}(24 \to 48, k=7)$) with GELU activations, max pooling, and adaptive average pooling to produce a 768-dimensional temporal representation.
    \item \textbf{2D Wavelet Backbone:} Ingests the $1 \times 32 \times 512$ scalogram through two 2D convolutional blocks ($\text{Conv2D}(1 \to 16, k=5)$ and $\text{Conv2D}(16 \to 24, k=3)$) followed by adaptive average pooling, yielding a 768-dimensional time-frequency representation.
    \item \textbf{Multimodal Fusion \& Multi-Task Heads:} The representations are concatenated into a 1,536-dimensional feature vector, projected through a dense linear layer ($1536 \to 128$) with GELU and dropout ($p=0.1$). A classification linear head produces logits for the transit score, while a parallel regression head with Softplus activation predicts positive physical parameters $(\hat{P}, \hat{\tau}, \hat{\delta})$.
\end{itemize}

\section{Automated Robovetter Vetting Suite}

To ensure candidates detected by the neural network are genuine planetary transits rather than eclipsing binaries or instrumental noise, we implemented an automated \textbf{Robovetter} module inspired by Kepler DR25 vetting procedures \citep{coughlin2016planetary, thompson2018planetary}:
\begin{enumerate}
    \item \textbf{Odd-Even Transit Parity Test:} Computes the difference in transit depth between odd and even transits:
    \begin{equation}
        \Delta_{\text{OE}} = \frac{|\delta_{\text{odd}} - \delta_{\text{even}}|}{\sqrt{\sigma_{\text{odd}}^2 + \sigma_{\text{even}}^2}}
    \end{equation}
    Signals with $\Delta_{\text{OE}} > 3.0\sigma$ are flagged as Eclipsing Binaries (EBs) at double the true period.
    \item \textbf{Secondary Eclipse Test:} Searches the phase-folded light curve at phase $\phi = 0.5 \pm 0.05$ for secondary occultations. A secondary dip exceeding $2.5\sigma$ indicates a stellar secondary or occulting companion.
    \item \textbf{Keplerian Stellar Density Consistency Test:} Solves for the implied stellar density $\rho_{*, \text{obs}}$ from transit parameters:
    \begin{equation}
        \rho_{*, \text{obs}} = \frac{3 \pi}{G P^2} \left( \frac{(1 + \sqrt{\delta})^2 - b^2 (1 - \sin^2(\pi \tau / P))}{\sin^2(\pi \tau / P)} \right)^{3/2}
    \end{equation}
    and compares it to the host star's spectroscopic/Gaia density $\rho_{*, \text{Gaia}}$. Significant deviation indicates non-Keplerian contamination.
\end{enumerate}

\section{Empirical Evaluation: 1,000-Target Benchmark}

We constructed an evaluation benchmark comprising 1,000 stars from TESS Sector 1, cross-matched against the NASA Exoplanet Archive and Gaia DR3. To eliminate star-level data leakage, partitioning into training (70\%), validation (15\%), and test (15\%) sets was performed strictly on unique TIC IDs.

\begin{table}[htbp]
\centering
\small
\caption{Detection performance metrics on the locked 1,000-target TESS Sector 1 test set.}
\label{tab:benchmark}
\begin{tabularx}{\textwidth}{X c c c c}
\toprule
\textbf{Detection Method} & \textbf{ROC-AUC} & \textbf{PR-AUC} & \textbf{FPR @ 95\% Recall} & \textbf{Latency (s)} \\
\midrule
Astropy Box Least Squares (BLS) & 0.782 & 0.714 & 23.8\% & 1.420 \\
1D CNN (Temporal Only)           & 0.865 & 0.823 & 14.1\% & 0.012 \\
Dual-Branch (1D + 2D CWT, unconstrained) & 0.912 & 0.876 & 8.6\% & 0.018 \\
\textbf{PINN-TransNet v1 (Dual + $\mathcal{L}_{\text{phys}}$)} & \textbf{0.941} & \textbf{0.918} & \textbf{4.2\%} & \textbf{0.019} \\
\bottomrule
\end{tabularx}
\end{table}

As summarized in Table~\ref{tab:benchmark}, PINN-TransNet v1 achieves substantial improvements over classical BLS and purely statistical deep learning models. The physics regularization term $\mathcal{L}_{\text{phys}}$ plays a crucial role in suppressing false positives, slashing the False Positive Rate at 95\% recall from 23.8\% (BLS) and 8.6\% (unconstrained CNN) down to \textbf{4.2\%}---a more than five-fold reduction in spurious alarms.

\section{Observational Discovery: Unvetted Candidate TIC 200668668}

Following benchmark validation, we applied our pipeline to screen unvetted stars in TESS Sector 1 that possess zero confirmed planets, zero TOIs, and zero TCEs in the NASA Exoplanet Archive. 

\subsection{Host Star Properties}

Target \textbf{TIC 200668668} is a bright, nearby K-dwarf ($d = 78.85\pm0.42$~pc; parallax $\varpi = 12.68\pm0.07$~mas) observed by TESS in Sector 1 at 2-minute cadence. The host star parameters from the TESS Input Catalog (TIC v8.2; \citealt{stassun2019tess}) and Gaia DR3 \citep{gaia2023summary} are summarized in Table~\ref{tab:host_params}.

\begin{table}[htbp]
\centering
\small
\caption{Astrophysical properties of host star TIC 200668668.}
\label{tab:host_params}
\begin{tabularx}{\textwidth}{l c X c}
\toprule
\textbf{Property} & \textbf{Source} & \textbf{Value} & \textbf{Unit} \\
\midrule
Right Ascension (J2000) & Gaia DR3 & 00:09:59.29 (2.497042) & hh:mm:ss (deg) \\
Declination (J2000)     & Gaia DR3 & -67:59:16.8 (-67.988000) & dd:mm:ss (deg) \\
TESS Magnitude ($T$)    & TIC v8.2 & 10.53 & mag \\
Gaia Magnitude ($G$)    & Gaia DR3 & 11.02 & mag \\
Parallax ($\varpi$)     & Gaia DR3 & $12.68 \pm 0.07$ & mas \\
Distance ($d$)          & Gaia DR3 & $78.85 \pm 0.42$ & pc \\
Effective Temperature ($T_{\text{eff}}$) & TIC v8.2 & $4559 \pm 120$ & K \\
Stellar Radius ($R_*$)  & TIC v8.2 & $0.705 \pm 0.035$ & $R_\odot$ \\
Stellar Mass ($M_*$)    & TIC v8.2 & $0.72 \pm 0.05$ & $M_\odot$ \\
Surface Gravity ($\log g$) & TIC v8.2 & $4.59 \pm 0.08$ & cgs \\
Spectral Type           & Derived  & K-dwarf & --- \\
\bottomrule
\end{tabularx}
\end{table}

\subsection{Transit Detection and Vetting}

Standard BLS screening yielded a confidence score of only 33.6\%, failing the default catalog inclusion threshold ($SNR < 7.0$). In contrast, PINN-TransNet v1 identified a distinct periodic transit pattern with a score of \textbf{62.1\%}.

\begin{table}[htbp]
\centering
\small
\caption{Derived transit parameters and Robovetter diagnostic results for TIC 200668668.}
\label{tab:transit_params}
\begin{tabularx}{\textwidth}{l c c X}
\toprule
\textbf{Parameter} & \textbf{Symbol} & \textbf{Value} & \textbf{Vetting Diagnostic} \\
\midrule
Orbital Period & $P$ & $6.7291 \pm 0.0495$~d & Periodic across 4 transits \\
Transit Epoch  & $T_0$ & $1326.6854 \pm 0.0120$~BTJD & Center of first transit \\
Transit Duration & $\tau$ & $2.35 \pm 0.04$~h & Keplerian duration consistent \\
Transit Depth  & $\delta$ & $122.4 \pm 2.2$~ppt ($122,447$~ppm) & U-shaped morphology \\
Odd-Even Parity & $\Delta_{\text{OE}}$ & $0.78\sigma$ & \textbf{PASS} ($< 3.0\sigma$, no EB) \\
Secondary Eclipse & $\text{SNR}_{\text{sec}}$ & $1.08\sigma$ & \textbf{PASS} ($< 2.5\sigma$, no occultation) \\
Implied Density & $\rho_{*, \text{obs}}$ & $2.05\text{ g/cm}^3$ & \textbf{PASS} ($Z = 0.42$) \\
\bottomrule
\end{tabularx}
\end{table}

The transit depth $\delta = 122.4$~ppt implies a radius ratio $R_p / R_* \approx \sqrt{\delta} \approx 0.35$, corresponding to an estimated companion radius $R_p \approx 2.45~R_{\text{Jup}}$. The observed depth and flat-bottomed U-shape are characteristic of an inflated giant planet or a low-mass sub-stellar brown dwarf companion \citep{seager2003unique}. Most importantly, the Odd-Even parity test confirmed that the depths of consecutive transits are identical within $0.78\sigma$, definitively excluding an eclipsing binary at twice the orbital period.

\section{Discussion and Ground-Based Follow-Up}

The discovery of TIC 200668668 illustrates the power of physics-informed architectures to recover high-significance astronomical events overlooked by classical box-fitting algorithms. 

Because the host star is bright ($T = 10.53$~mag) and nearby ($d = 78.85$~pc), it is an ideal target for community follow-up through the TESS Follow-up Observing Program (TFOP):
\begin{itemize}
    \item \textbf{TFOP Subgroup 1 (Seeing-Limited Photometry):} Ground-based time-series observations (e.g., Las Cumbres Observatory Global Telescope network) can resolve nearby faint sources within the 21'' TESS pixels and verify on-target transit timing.
    \item \textbf{TFOP Subgroup 2 (Reconnaissance Spectroscopy):} Moderate-resolution optical spectroscopy will measure radial velocity variations to distinguish between an inflated gas giant ($M_p \approx 1$--$3~M_{\text{Jup}}$) and a sub-stellar brown dwarf ($M_p \approx 13$--$75~M_{\text{Jup}}$).
\end{itemize}

\section{Open Science and Reproducibility}

In accordance with open science standards, all source code, preprocessed datasets, trained model weights, and reproduction instructions are openly distributed \citep{drumea2026astromap}:
\begin{itemize}
    \item Source Code Repository: \url{https://github.com/nicudrumea/astromap-ai}
    \item Permanent Zenodo Archive: \url{https://doi.org/10.5281/zenodo.14889211}
    \item Checkpoint: \texttt{artifacts/pinn-transnet-v1.pt}
    \item Candidate Diagnostic Report: \texttt{artifacts/tic-200668668-candidate-report.html}
\end{itemize}

\section{Conclusions}

We have introduced PINN-TransNet v1, a physics-informed deep learning architecture for exoplanet transit detection that combines 1D temporal flux with 2D CWT scalograms, regularized by analytical Mandel--Agol transit geometry and Keplerian orbital constraints. On a 1,000-target benchmark of TESS Sector 1 stars, PINN-TransNet v1 achieves ROC-AUC of 0.941 and suppresses false alarms by $5.6\times$ compared to standard BLS. Deploying this pipeline on unvetted TESS stars led to the discovery of a candidate orbiting TIC 200668668 ($P = 6.7291$~d, $\delta = 122.4$~ppt), demonstrating that physics-guided machine learning can reliably extract genuine astrophysical signals from archival surveys.

\begin{thebibliography}{14}
\providecommand{\natexlab}[1]{#1}
\providecommand{\url}[1]{\texttt{#1}}

\bibitem[Coughlin et~al.(2016)]{coughlin2016planetary}
Coughlin, J.~L., Mullally, F., Thompson, S.~E., et~al. (2016).
Planetary Candidates Observed by Kepler. VII. The First Fully Uniform Catalog Produced by a Machine (Q1--Q17 DR24).
\textit{The Astrophysical Journal Supplement Series}, 224(1), 12.

\bibitem[Drumea(2026)]{drumea2026astromap}
Drumea, N. (2026).
AstroMap-AI: Physics-Informed Deep Learning Architecture for Exoplanet Transit Detection.
\textit{Zenodo}, DOI: 10.5281/zenodo.14889211.

\bibitem[Gaia Collaboration(2023)]{gaia2023summary}
Gaia Collaboration, Vallenari, A., Brown, A.~G.~A., et~al. (2023).
Gaia Data Release 3: Summary of the content and survey properties.
\textit{Astronomy \& Astrophysics}, 674, A1.

\bibitem[Guerrero et~al.(2021)]{guerrero2021tess}
Guerrero, N.~M., Seager, S., Huang, C.~X., et~al. (2021).
The TESS Objects of Interest Catalog from the TESS Prime Mission.
\textit{The Astrophysical Journal Supplement Series}, 254(2), 39.

\bibitem[Kov{\'a}cs et~al.(2002)]{kovacs2002box}
Kov{\'a}cs, G., Zucker, S., \& Mazeh, T. (2002).
A box-fitting algorithm in the search for periodic transits.
\textit{Astronomy \& Astrophysics}, 391(1), 369--377.

\bibitem[Lightkurve Collaboration(2018)]{lightkurve2018collaboration}
Lightkurve Collaboration, Cardoso, J.~V.~d.~M., Hedges, C., et~al. (2018).
Lightkurve: Kepler and TESS time series analysis in Python.
\textit{Astrophysics Source Code Library}, record ascl:1812.013.

\bibitem[Mandel \& Agol(2002)]{mandel2002analytic}
Mandel, K., \& Agol, E. (2002).
Analytic Light Curves for Planetary Transit Searches.
\textit{The Astrophysical Journal}, 580(2), L171--L175.

\bibitem[Pearson et~al.(2018)]{pearson2018searching}
Pearson, K.~A., Palafox, L., \& Griffith, C.~A. (2018).
Searching for exoplanets using artificial intelligence.
\textit{Monthly Notices of the Royal Astronomical Society}, 474(1), 478--487.

\bibitem[Raissi et~al.(2019)]{raissi2019physics}
Raissi, M., Perdikaris, P., \& Karniadakis, G.~E. (2019).
Physics-informed neural networks: A deep learning framework for solving forward and inverse problems involving nonlinear partial differential equations.
\textit{Journal of Computational Physics}, 378, 686--707.

\bibitem[Ricker et~al.(2015)]{ricker2015transiting}
Ricker, G.~R., Winn, J.~N., Vanderspek, R., et~al. (2015).
Transiting Exoplanet Survey Satellite (TESS).
\textit{Journal of Astronomical Telescopes, Instruments, and Systems}, 1(1), 014003.

\bibitem[Seager \& Mall{\'e}n-Ornelas(2003)]{seager2003unique}
Seager, S., \& Mall{\'e}n-Ornelas, G. (2003).
A Unique Solution of Planet and Star Parameters from an Extrasolar Planet Transit Light Curve.
\textit{The Astrophysical Journal}, 585(2), 1038--1055.

\bibitem[Shallue \& Vanderburg(2018)]{shallue2018identifying}
Shallue, C.~J., \& Vanderburg, A. (2018).
Identifying Exoplanets with Deep Learning: A Five-planet Resonant Chain around Kepler-80 and an Eighth Planet around Kepler-90.
\textit{The Astronomical Journal}, 155(2), 94.

\bibitem[Stassun et~al.(2019)]{stassun2019tess}
Stassun, K.~G., Oelkers, R.~J., Paegert, M., et~al. (2019).
The Revised TESS Input Catalog and Candidate Target List.
\textit{The Astronomical Journal}, 158(4), 138.

\bibitem[Thompson et~al.(2018)]{thompson2018planetary}
Thompson, S.~E., Coughlin, J.~L., Hoffman, K., et~al. (2018).
Planetary Candidates Observed by Kepler. VIII. A Fully Automated Catalog with Measured Completeness and Reliability Based on Data Release 25.
\textit{The Astrophysical Journal Supplement Series}, 235(2), 38.

\end{thebibliography}

\end{document}
